\documentclass[12pt, a4paper]{article}

% ── Кодировка и язык ──────────────────────────────────────────────────────────
\usepackage[T2A]{fontenc}
\usepackage[utf8]{inputenc}
\usepackage[russian]{babel}

% ── Геометрия страницы ────────────────────────────────────────────────────────
\usepackage[
  top=20mm, bottom=25mm,
  left=20mm, right=20mm
]{geometry}

% ── Математика ────────────────────────────────────────────────────────────────
\usepackage{amsmath, amssymb}

% ── Графика и таблицы ─────────────────────────────────────────────────────────
\usepackage{graphicx}
\usepackage{booktabs}
\usepackage{array}
\usepackage{caption}
\usepackage{subcaption}
\usepackage{float}

% ── Прочее оформление ─────────────────────────────────────────────────────────
\usepackage{hyperref}
\usepackage{xcolor}
\usepackage{enumitem}
\usepackage{microtype}
\usepackage{indentfirst}
\usepackage{setspace}
\usepackage{parskip}

\setlength{\parindent}{1.25cm}
\setlength{\parskip}{0pt}
\onehalfspacing

\hypersetup{
  colorlinks = true,
  linkcolor  = black,
  citecolor  = black,
  urlcolor   = blue
}

% ──────────────────────────────────────────────────────────────────────────────
\begin{document}

% ── Титульная страница ────────────────────────────────────────────────────────
\begin{titlepage}
  \centering
  \vspace*{2cm}

  {\Large\bfseries Классификация рамановских спектров\\
    мозговой ткани методами машинного обучения}\par
  \vspace{0.5cm}
  {\large Команда: Gradient Island}\par

  \vspace{3cm}

  \begin{tabular}{rl}
    \textbf{Задача:}       & Nuclear It Hack \\[4pt]
    \textbf{Метод:}        & Ансамблевое машинное обучение со стекингом \\[4pt]
    \textbf{Метрика:}      & Macro F1-score \\[4pt]
    \textbf{Результат:}    & Macro F1 = 0.7908 (OOF LOAO-CV) \\
  \end{tabular}

  \vfill
  {\small \today}
\end{titlepage}

% ── Оглавление ────────────────────────────────────────────────────────────────
\tableofcontents
\newpage

% ══════════════════════════════════════════════════════════════════════════════
\section{Постановка задачи}
% ══════════════════════════════════════════════════════════════════════════════

Целью работы является разработка модели машинного обучения, которая
по рамановскому спектру мозговой ткани крысы определяет
экспериментальную группу животного: \textit{контроль},
\textit{эндогенная экспрессия HSP70} (endo) или
\textit{экзогенная экспрессия HSP70} (exo).
Решаемая задача мотивирована необходимостью выявления спектральных
маркеров, наиболее чувствительных к нейропротективным
и стресс-индуцированным молекулярным изменениям в ткани мозга.

Рамановская спектроскопия позволяет неинвазивно регистрировать
молекулярный состав ткани через характеристические полосы
колебаний химических связей. Белок теплового шока HSP70 является
ключевым шапероном клеточного стресс-ответа: его эндогенная
индукция возникает в ответ на повреждение, тогда как экзогенное
введение рассматривается как потенциальная нейропротективная стратегия.
Спектральные различия между тремя группами ожидаются в диапазонах,
связанных с белковой конформацией (амидные полосы), липидным составом
мембран и нуклеиновыми кислотами.

\subsection{Формальная постановка}

Пусть имеется набор данных
$\mathcal{D} = \{(\mathbf{x}_i,\, y_i)\}_{i=1}^{N}$,
где $\mathbf{x}_i \in \mathbb{R}^p$ — вектор признаков,
извлечённых из рамановского спектра $i$-й спектральной карты,
а $y_i \in \{\text{control},\, \text{endo},\, \text{exo}\}$ —
метка класса.
Требуется построить классификатор
$f\colon \mathbb{R}^p \to \{\text{control}, \text{endo}, \text{exo}\}$,
максимизирующий макро-усреднённую $F_1$-меру:

\begin{equation}
  F_1^{\text{macro}} = \frac{1}{3}\sum_{c \in \mathcal{C}} F_1^{(c)},
  \qquad
  F_1^{(c)} = \frac{2\,\mathrm{TP}_c}{2\,\mathrm{TP}_c + \mathrm{FP}_c + \mathrm{FN}_c}.
  \label{eq:f1macro}
\end{equation}

\subsection{Данные}

Исходные данные представлены спектральными картами различных
отделов мозга крыс, полученных методом конфокальной рамановской
микроскопии. Спектры регистрировались в двух спектральных окнах:

\begin{itemize}[leftmargin=1.5cm]
  \item \textbf{Диапазон 1500~см$^{-1}$} (600–1800~см$^{-1}$) —
        «отпечаток пальца» (fingerprint region): белковые, нуклеиновые
        и липидные полосы поглощения.
  \item \textbf{Диапазон 2900~см$^{-1}$} (2800–3100~см$^{-1}$) —
        валентные колебания C–H: характеристика липидного состава мембран.
\end{itemize}

Суммарный объём: 118–119 спектральных карт на диапазон,
4 животных (по числу LOAO-фолдов), три класса
(40 control, 36–37 endo, 42 exo).
Исходный датафрейм содержит порядка 126 миллионов строк вида
\texttt{(X, Y, Wave, Intensity, label, animal, center, brain, place)}.


% ══════════════════════════════════════════════════════════════════════════════
\section{Методология}
% ══════════════════════════════════════════════════════════════════════════════

% ──────────────────────────────────────────────────────────────────────────────
\subsection{Предобработка спектров}
% ──────────────────────────────────────────────────────────────────────────────

\paragraph{Усреднение и построение общей сетки.}
Для каждой спектральной карты (уникальная комбинация животного,
отдела мозга и места измерения) точки одного пространственного
положения $(X, Y)$ усреднялись по интенсивности при одинаковом
значении сдвига. Тем самым каждая карта редуцировалась к единому
среднему спектру. Затем все спектры интерполировались на
унифицированную сетку из 1500 равноотстоящих точек в диапазоне
$[q_{2\%},\, q_{98\%}]$ по оси волновых чисел.

\paragraph{Коррекция флуоресцентного фона.}
Флуоресценция тканей создаёт широкую аддитивную базовую линию,
маскирующую рамановский сигнал. Для её удаления применялся
алгоритм AsLS (\textit{Asymmetric Least Squares Smoothing}):

\begin{equation}
  \hat{\mathbf{z}} = \arg\min_{\mathbf{z}}
    \Bigl[
      \mathbf{w}^\top (\mathbf{y} - \mathbf{z})^2
      + \lambda \|\mathbf{D}^{(2)}\mathbf{z}\|^2
    \Bigr],
  \label{eq:als}
\end{equation}

где $\mathbf{D}^{(2)}$ — матрица вторых разностей,
$\lambda = 10^5$ управляет гладкостью базовой линии,
веса $w_j$ итеративно обновляются: $w_j = p$ при $y_j > z_j$ и
$w_j = 1-p$ иначе ($p = 0.01$, число итераций 15).
Скорректированный спектр: $\tilde{\mathbf{y}} = \mathbf{y} - \hat{\mathbf{z}}$.

\paragraph{Сглаживание.}
К скорректированным спектрам применялся фильтр
Савицкого–Голая (\textit{Savitzky-Golay}): ширина окна 11 точек,
степень полинома 3. Данный фильтр сохраняет форму пиков,
уменьшая высокочастотный шум.

\paragraph{Стандартное нормальное варьирование (SNV).}
Для устранения мультипликативных рассеивающих эффектов
каждый спектр нормализовался по формуле:

\begin{equation}
  x_j^{\text{snv}} = \frac{x_j - \bar{x}}{\sigma_x},
  \label{eq:snv}
\end{equation}

где $\bar{x}$ и $\sigma_x$ — среднее и стандартное отклонение
интенсивностей данного спектра.

Результат каждого шага предобработки иллюстрирует рис.~\ref{fig:preproc}.

\begin{figure}[H]
  \centering
  \includegraphics[width=1.0\textwidth]{../notebooks/plots/1500_baseline_overlay.png}
\end{figure}

% ──────────────────────────────────────────────────────────────────────────────
\subsection{Извлечение признаков}
% ──────────────────────────────────────────────────────────────────────────────

Признаковое пространство формировалось из нескольких групп,
каждая из которых кодирует различный аспект спектра.

\paragraph{Признаки биохимических полос.}
Для каждого из диагностически значимых диапазонов (табл.~\ref{tab:bands})
вычислялись: площадь под кривой (численное интегрирование методом трапеций),
максимальная и средняя интенсивность, положение главного пика,
стандартное отклонение и коэффициент асимметрии. Дополнительно
формировались попарные отношения площадей соседних полос.

\begin{table}[H]
\centering
\caption{Биохимически значимые рамановские полосы.}
\label{tab:bands}
\begin{tabular}{llr}
\toprule
\textbf{Диапазон} & \textbf{Молекулярное отнесение} & \textbf{Позиция, см$^{-1}$} \\
\midrule
\multicolumn{3}{l}{\textit{Fingerprint region (1500)}} \\
& Фенилаланин                & 1000–1010 \\
& Пролин                     & 850–860   \\
& Гидроксипролин             & 875–885   \\
& Амид III (\(\beta\)-лист)          & 1230–1270 \\
& Амид III (\(\alpha\)-спираль)       & 1270–1310 \\
& Деформация CH$_2$          & 1440–1470 \\
& Амид II                    & 1540–1560 \\
& Амид I (\(\alpha\)-спираль)         & 1650–1660 \\
& Амид I (\(\beta\)-лист)            & 1670–1690 \\
& Эфир липидов               & 1730–1760 \\
& ДНК/РНК                    & 780–800   \\
& Тирозин                    & 830–840   \\
\midrule
\multicolumn{3}{l}{\textit{CH-region (2900)}} \\
& CH$_2$ симметр. валентные  & 2845–2855 \\
& CH$_2$ асимметр. валентные & 2880–2900 \\
& CH$_3$ валентные           & 2930–2960 \\
& Олефиновые CH              & 3000–3030 \\
\bottomrule
\end{tabular}
\end{table}

\paragraph{Производные признаки.}
Первая и вторая производные SNV-спектра (метод Савицкого–Голая)
характеризуют скорость и кривизну изменения сигнала.
Для каждой производной вычислялись: среднее, стандартное отклонение,
максимум абсолютного значения, суммарная энергия,
число нулевых переходов, зональные энергии (8 равных сегментов).

\paragraph{Вейвлет-признаки.}
Дискретное вейвлет-преобразование (\textit{Daubechies db8}, 5 уровней)
разлагало спектр на поддиапазоны различных масштабов.
Для каждого коэффициентного вектора извлекались: суммарная энергия,
относительная энергия, среднее, стандартное отклонение и вейвлет-энтропия
$H = -\sum_k p_k \ln p_k$, где $p_k = c_k^2 / \sum_j c_j^2$.

\paragraph{Зональные статистики.}
Спектр делился на 8 равных частей; для каждой зоны вычислялись
среднее и стандартное отклонение интенсивности, отношения к первой зоне
и разности между соседними зонами.

\paragraph{Снижение размерности (PCA и NMF).}
Из SNV-матрицы карт извлекались 20 главных компонент (PCA)
и 15 неотрицательных компонент (NMF, предварительная Min-Max нормализация).
Обе трансформации фиксировались при обучении и применялись
без переобучения при инференсе.

После конкатенации всех блоков итоговое признаковое пространство
составляет 284–328 признаков на диапазон.

\paragraph{Отбор признаков.}
Для каждого диапазона применялся LightGBM-ранкинг по важности
типа \textit{gain}: из полного набора отбиралось 120 признаков
с наибольшим вкладом. Это позволяло снизить размерность без потери
информативности (рис.~\ref{fig:importance}).

\begin{figure}[H]
  \centering
  \fbox{\parbox{0.9\textwidth}{\centering\vspace{3cm}
    \textit{[Место для графика важности признаков по gain (top-30
    для каждого из диапазонов 1500 и 2900)]}
  \vspace{3cm}}}
  \caption{Топ-30 признаков по gain-важности LightGBM для диапазонов
    1500~см$^{-1}$ (слева) и 2900~см$^{-1}$ (справа).}
  \label{fig:importance}
\end{figure}

% ──────────────────────────────────────────────────────────────────────────────
\subsection{Стратегия валидации: Leave-One-Animal-Out}
% ──────────────────────────────────────────────────────────────────────────────

Принципиальной проблемой оценки модели на биологических данных
является утечка информации между образцами одного животного:
спектры с разных участков мозга одной крысы биологически зависимы.
Использование обычной $k$-кратной перекрёстной проверки
приводит к завышенным оценкам.

Для корректной оценки обобщающей способности применялась схема
\textbf{Leave-One-Animal-Out (LOAO)}: на каждой итерации все карты
одного животного ($n = 4$ животных) изымались в тест,
остальные использовались для обучения (табл.~\ref{tab:loao}).
Все метрики (включая настройку гиперпараметров Optuna)
вычислялись исключительно на out-of-fold предсказаниях.

\begin{table}[H]
\centering
\caption{Структура LOAO-фолдов.}
\label{tab:loao}
\begin{tabular}{ccc}
\toprule
\textbf{Фолд} & \textbf{Животное в тесте} & \textbf{Карт в тесте} \\
\midrule
1 & 1 & 36 \\
2 & 2 & 18 \\
3 & 3 & 31 \\
4 & 4 & 33 \\
\bottomrule
\end{tabular}
\end{table}

% ──────────────────────────────────────────────────────────────────────────────
\subsection{Базовые модели и подбор гиперпараметров}
% ──────────────────────────────────────────────────────────────────────────────

Обучались семь базовых классификаторов:
XGBoost, LightGBM, CatBoost, Random Forest, Extra Trees, SVM (RBF/poly),
MLP. Перед подачей в каждую модель признаки масштабировались
\texttt{StandardScaler}.

Гиперпараметры подбирались методом байесовской оптимизации
(библиотека Optuna, семплер TPE, 50 испытаний на модель).
Целевой функцией служил macro $F_1$ по LOAO-OOF предсказаниям.

Каждая модель обучалась независимо на трёх наборах признаков:
диапазон 1500, диапазон 2900 и \textit{joint} (конкатенация обоих
диапазонов с дополнительными межканальными признаками).
Итого обучалось $7 \times 3 = 21$ базовая модель.

Результаты базовых моделей приведены на рис.~\ref{fig:base_models}.

\begin{figure}[H]
  \centering
  \fbox{\parbox{0.9\textwidth}{\centering\vspace{3cm}
    \textit{[Место для столбчатой диаграммы macro F1 (OOF)
    всех 21 базовой модели]}
  \vspace{3cm}}}
  \caption{Macro F1 (LOAO-OOF) для каждой из 21 базовой модели,
    сгруппированных по диапазону.}
  \label{fig:base_models}
\end{figure}

% ──────────────────────────────────────────────────────────────────────────────
\subsection{Ансамблирование: стекинг с оптимизацией порогов}
% ──────────────────────────────────────────────────────────────────────────────

OOF-вероятности всех 21 модели конкатенировались в мета-матрицу
$\mathbf{M} \in \mathbb{R}^{N \times (21 \cdot 3)}$ и подавались
на вход мета-классификатору. Рассматривались следующие стратегии:

\begin{enumerate}[leftmargin=1.5cm]
  \item \textbf{Soft Voting} — взвешенное среднее вероятностей,
        где веса пропорциональны OOF $F_1$ каждой модели.
  \item \textbf{Geometric Mean} — геометрическое среднее вероятностей,
        сглаживающее эффект уверенных ошибочных предсказаний.
  \item \textbf{Rank Averaging} — усреднение ранговых позиций
        по каждому классу.
  \item \textbf{Stacking} — мета-классификатор (Logistic Regression),
        обученный на $\mathbf{M}$ с 5-кратной стратифицированной
        перекрёстной проверкой.
  \item \textbf{Best + Threshold} — лучшая стратегия с дополнительной
        покоординатной оптимизацией порогов решения:
        $\hat{y} = \arg\max_c (p_c / \theta_c)$,
        где $\theta_c$ подбирается сеточным поиском по
        $F_1^{\text{macro}}$.
\end{enumerate}

Итоговые OOF-метрики всех стратегий представлены в
табл.~\ref{tab:ensemble}.

\begin{table}[H]
\centering
\caption{Сравнение стратегий ансамблирования (LOAO-OOF).}
\label{tab:ensemble}
\begin{tabular}{lcccc}
\toprule
\textbf{Стратегия} & \textbf{Macro F1} & \textbf{Accuracy} &
\textbf{F1\,control} & \textbf{F1\,endo} \\
\midrule
Soft Voting      & 0.494 & 0.500 & 0.490 & 0.390 \\
Geometric Mean   & 0.504 & —     & —     & —     \\
Rank Averaging   & 0.486 & —     & —     & —     \\
Stacking (LR)    & 0.764 & 0.763 & 0.750 & 0.810 \\
\textbf{Best + Threshold} & \textbf{0.791} & \textbf{0.788} &
\textbf{0.770} & \textbf{0.850} \\
\bottomrule
\end{tabular}
\end{table}


% ══════════════════════════════════════════════════════════════════════════════
\section{Результаты}
% ══════════════════════════════════════════════════════════════════════════════

% ──────────────────────────────────────────────────────────────────────────────
\subsection{Итоговые метрики классификации}
% ──────────────────────────────────────────────────────────────────────────────

Лучший результат достигнут стратегией «Stacking + Threshold»:
macro $F_1 = 0.791$, accuracy $= 0.788$ на LOAO-OOF.
Детальный отчёт классификации представлен в табл.~\ref{tab:clf_report}.

\begin{table}[H]
\centering
\caption{Отчёт классификации (LOAO-OOF, Best+Threshold).}
\label{tab:clf_report}
\begin{tabular}{lcccc}
\toprule
\textbf{Класс} & \textbf{Precision} & \textbf{Recall} &
\textbf{F1-score} & \textbf{Support} \\
\midrule
control  & 0.76 & 0.78 & 0.77 & 40 \\
endo     & 0.86 & 0.83 & 0.85 & 36 \\
exo      & 0.76 & 0.76 & 0.76 & 42 \\
\midrule
macro avg & 0.79 & 0.79 & 0.79 & 118 \\
\bottomrule
\end{tabular}
\end{table}

Матрицы ошибок для всех стратегий приведены на рис.~\ref{fig:cm}.

\begin{figure}[H]
  \centering
  \fbox{\parbox{0.9\textwidth}{\centering\vspace{3.5cm}
    \textit{[Место для матриц ошибок (confusion matrices) по пяти
    стратегиям ансамблирования]}
  \vspace{3.5cm}}}
  \caption{Матрицы ошибок LOAO-OOF для пяти стратегий ансамблирования.
    Слева направо: Soft Voting, Geometric Mean, Rank Averaging,
    Stacking, Best+Threshold.}
  \label{fig:cm}
\end{figure}

% ──────────────────────────────────────────────────────────────────────────────
\subsection{Анализ значимости признаков}
% ──────────────────────────────────────────────────────────────────────────────

Наиболее высокий gain-вклад в обоих диапазонах вносят:

\begin{itemize}[leftmargin=1.5cm]
  \item \textbf{1500~см$^{-1}$:} отношения интегралов амид~I / амид~III
        (маркер вторичной структуры белка), площадь полосы эфиров
        липидов (1730–1760~см$^{-1}$), вейвлет-энергии среднего
        масштаба ($cD_3$, $cD_4$).
  \item \textbf{2900~см$^{-1}$:} отношение CH$_2$/CH$_3$
        (степень ненасыщенности жирных кислот),
        площадь олефиновой полосы (3000–3030~см$^{-1}$),
        PCA-компонента~1 (доля объяснённой дисперсии $\approx 40\%$).
\end{itemize}

Это согласуется с биологическими ожиданиями: экспрессия HSP70
сопровождается реорганизацией белков и изменениями
в мембранных липидах нейронов.

Средние нормализованные спектры по классам представлены на
рис.~\ref{fig:mean_spectra}.

\begin{figure}[H]
  \centering
  \fbox{\parbox{0.9\textwidth}{\centering\vspace{3.5cm}
    \textit{[Место для графиков средних SNV-спектров по трём классам
    (control / endo / exo) для диапазонов 1500 и 2900 с отмеченными
    значимыми полосами]}
  \vspace{3.5cm}}}
  \caption{Средние SNV-спектры трёх классов. Вертикальными пунктирными
    линиями отмечены биохимически значимые полосы с наибольшим
    вкладом в модель.}
  \label{fig:mean_spectra}
\end{figure}

% ──────────────────────────────────────────────────────────────────────────────
\subsection{Пространство признаков: визуализация PCA}
% ──────────────────────────────────────────────────────────────────────────────

\begin{figure}[H]
  \centering
  \fbox{\parbox{0.9\textwidth}{\centering\vspace{3.5cm}
    \textit{[Место для scatter-графика PC1 vs PC2 с раскраской по классам
    для диапазонов 1500 и 2900]}
  \vspace{3.5cm}}}
  \caption{Проекция спектральных карт на плоскость первых двух
    главных компонент (PCA, SNV-спектры). Слева: диапазон
    1500~см$^{-1}$; справа: 2900~см$^{-1}$.}
  \label{fig:pca_scatter}
\end{figure}


% ══════════════════════════════════════════════════════════════════════════════
\section{Обсуждение}
% ══════════════════════════════════════════════════════════════════════════════

Достигнутый macro $F_1 = 0.791$ демонстрирует, что рамановские
спектры мозговой ткани несут устойчивую молекулярную «подпись»
экспериментальной группы, различимую автоматическими методами.

Наблюдаемое превосходство стекинга над прямыми стратегиями агрегации
(Soft Voting: 0.494 против 0.764) объясняется тем, что мета-классификатор
обучается учитывать различные систематические смещения базовых моделей.
Ни одна базовая модель в отдельности не превышала macro $F_1 = 0.55$,
что указывает на принципиальную роль ансамблирования.

Класс \textit{endo} демонстрирует наилучшую разделимость ($F_1 = 0.85$),
что соответствует биологической гипотезе: эндогенная индукция HSP70
является сильным стрессовым ответом, порождающим выраженные
спектральные изменения (денатурация белков, окислительные модификации
липидов). Смешение \textit{control} и \textit{exo} несколько выше,
что логично: введение экзогенного HSP70 частично нормализует
молекулярный фенотип ткани.

Четыре LOAO-фолда — ограниченный объём для надёжной оценки
обобщающей способности. Расширение выборки новыми животными
является приоритетной задачей.

% ══════════════════════════════════════════════════════════════════════════════
\section{Заключение}
% ══════════════════════════════════════════════════════════════════════════════

В работе разработан полный пайплайн классификации рамановских спектров
мозговой ткани крысы по трём экспериментальным группам.
Ключевые результаты:

\begin{enumerate}[leftmargin=1.5cm]
  \item Разработана процедура предобработки, включающая коррекцию
        флуоресцентного фона (ALS), сглаживание (SG) и SNV-нормализацию,
        обеспечивающая воспроизводимое сравнение спектров разных карт.
  \item Сформировано многоуровневое признаковое пространство из биохимических
        полос, производных, вейвлет-коэффициентов и компонент
        снижения размерности.
  \item Применена схема LOAO-CV, исключающая утечку биологической
        информации между животными.
  \item Обучен ансамбль из 21 базовой модели (7 алгоритмов × 3 набора
        признаков) с байесовской настройкой гиперпараметров.
  \item Итоговый macro $F_1 = 0.791$ достигнут через стекинговый
        мета-классификатор с оптимизацией порогов решения.
\end{enumerate}

Наиболее информативными оказались признаки амидных полос
(вторичная структура белка) и соотношение CH$_2$/CH$_3$
(ненасыщенность мембранных липидов), что открывает перспективы
для разработки быстрых спектральных биомаркеров нейропротекции.

\end{document}
